%! TEX root = main.tex
\begin{solution}{ex:nt-1}
  Using the product formula:
  \[
    \begin{gathered}
      \varphi (72) = 72 \prod_{p\mid 72} (1-\frac{1}{p}), \qquad (p=2,3)\\
      \varphi (72) = 72 \cdot (\frac{1}{2})(\frac{2}{3})\\
      \varphi (72) = 72 \cdot \frac{1}{3}\\
      \varphi (72) = 24
    \end{gathered}
  \]
  The residues coprime to 72 are:\\ $1,5,7,11,13,17,19,23,25,29,31,35,37,41,43,47,49,53,55,59,61,65,67,71$\\
\end{solution}
\begin{solution}{ex:nt-2}
  Using Euler's algorithm and then back-substituting afterwards:
  \[
    \begin{gathered}
      64=3\cdot 17 + 13\\
      17=1\cdot 13 + 4\\
      13=3\cdot 4 + 1\\
    \end{gathered}
  \]
  64 and 17 are coprime. Now back-substituting:
  \[
    \begin{aligned}
      1&=13-3\cdot4\\
      1&=13-3(17-13)\\
      1&=4\cdot 13-3\cdot 17\\
      1&=4(64-3\cdot 17)-3\cdot 17\\
      1&=4\cdot 64 -15 \cdot 17\\
    \end{aligned}
  \]
  The Bézout coefficients are thus 4 and -15.
  This means that:
  \[
    \begin{gathered}
      -15\cdot17 \equiv 1 \quad\text{(mod 64)}\\
      49 \cdot 17 \equiv 1\quad \text{(mod 64)}
    \end{gathered}
  \]
  So the multiplicative inverse of 17 mod 64 is 49.

\end{solution}
\begin{solution}{ex:nt-3}
  Since 7 is prime and $3\nmid 7$, using Fermat's little theorem, we have:
  \[
    \begin{gathered}
      3^6\equiv 1 \quad \text{(mod 7)}\\
      3^{100}=(3^6)^{16}\cdot 3^4\\
      (3^6)^{16}\cdot 3^4\equiv 3^4 \quad \text{(mod 7)}\\
      3^4=81\equiv 4 \quad \text{(mod 7)}\\
    \end{gathered}
  \]
  Thus, we have:
  \[
    3^{100}\equiv 4 \quad \text{(mod 7)}
  \]
  
\end{solution}
\begin{solution}{ex:nt-4}
  Since 3 and 100 are coprime, we can use this identity:
  \[
    \begin{gathered}
      3^{\varphi (100)}\equiv 1 \quad \text{(mod 100)}\\
      3^{40} \equiv 1 \quad \text{(mod 100)}\\
      3^{2026}=(3^{40})^{50}\cdot 3^{26}\equiv 3^{26} \quad \text{(mod 100)}\\
    \end{gathered}
  \]
  Now, let's use CRT, with $M=100$. We can set $m_1=4$ and $m_2=25$ and find the system of equations:
  \[
    \begin{aligned}
      3^{\varphi(4)}&\equiv 1 \quad \text{(mod 4)}\\
      3^{2}&\equiv 1 \quad \text{(mod 4)}\\
      (3^{2})^{13} &= 3^{26}\\
      3^{26}&\equiv 1 \quad \text{(mod 4)}\\\\
      3^{\varphi (25)} &\equiv 1 \quad \text{(mod 25)}\\
      3^{20}&\equiv 1 \quad \text{(mod 25)}\\
      3^{26}&=3^{20}\cdot 3^6 \\
      3^{26}&\equiv 3^6 \equiv 4 \quad \text{(mod 25)}
    \end{aligned}
  \]
  Now, we have $M_1=25$ and $M_2=4$. Inverses: $25\equiv1\quad \text{(mod 4)}$ with $e_1=1$, and $4\equiv 4\quad \text{(mod 25)}$ with $e_2=19$. Now for the final step:
  \[
    \begin{gathered}
      3^{26}\equiv \sum_i a_i M_i e_i \quad \text{(mod 100)}\\
      3^{26}\equiv 1\cdot 25 \cdot 1 + 4\cdot 4 \cdot 19 \quad \text{(mod 100)}\\
      3^{26}\equiv 25 + 304 \quad \text{(mod 100)}\\
      3^{26} \equiv 329 \quad \text{(mod 100)}\\
      3^{26} \equiv 29 \quad \text{(mod 100)}
    \end{gathered}
  \]
  So the last 2 digits of $3^{2026}$ are 2 and 9.
\end{solution}
\begin{solution}{ex:nt-5}
  Using CRT with $M=4\cdot 5\cdot 7=140$ and $M_1=35$, $M_2=28$, $M_3=20$. Finding the inverses:
  \[
    \begin{gathered}
      35\equiv 3 \quad \text{(mod 4)}\\
      3\cdot 35 \equiv 1 \quad \text{(mod 4)}\\\\
      28\equiv 3 \quad \text{(mod 5)}\\
      2\cdot 28 \equiv 1 \quad \text{(mod 5)}\\\\
      20\equiv 6 \quad \text{(mod 7)}\\
      6\cdot 20 \equiv 1 \quad \text{(mod 7)}\\
    \end{gathered}
  \]
  Now let's use CRT to find the answer:
  \[
    \begin{gathered}
      x\equiv \sum_i a_i M_i e_i \quad \text{(mod M)}\\
      x \equiv 1 \cdot 35 \cdot 3 + 2\cdot 28 \cdot 2 + 3\cdot 20 \cdot 6 \quad \text{(mod M)}\\
      x \equiv 105 + 112 + 360 \quad \text{(mod M)}\\
      x \equiv 577 (mod M)\\
      x \equiv 17 \quad \text{(mod M)}
    \end{gathered}
  \]
  The smallest non-negative solution is 17, and all of the solutions are of the form $nM+17, n\in \N$
\end{solution}
